\documentclass[11pt]{amsart}

\usepackage[T1]{fontenc}
\usepackage{lmodern}
\usepackage{microtype}
\usepackage{amsmath,amssymb,amsthm}
\usepackage{booktabs}
\usepackage[hidelinks]{hyperref}

\hypersetup{
  pdftitle={A Negative Answer to Mohammadi's Question on Powers of Cover Ideals},
  pdfsubject={A counterexample concerning ordinary powers of graph cover ideals},
  pdfkeywords={cover ideal, linear resolution, Cohen--Macaulay graph, upper Koszul complex}
}

\newtheorem{theorem}{Theorem}
\theoremstyle{remark}
\newtheorem{remark}[theorem]{Remark}

\newcommand{\Ind}{\operatorname{Ind}}

\title[A nonlinear square of a cover ideal]
  {A Negative Answer to Mohammadi's Question on Powers of Cover Ideals}
\date{}
\subjclass[2020]{13D02, 13F55, 05E40}
\keywords{cover ideal, linear resolution, Cohen--Macaulay graph, upper Koszul complex}

\begin{document}

\begin{abstract}
We exhibit an eight-vertex Cohen--Macaulay graph $G$ whose cover
ideal $J(G)$ has a $5$-linear resolution over every field, while its
ordinary square does not have a linear resolution.  In fact,
\[
\beta_{1,(1,1,2,2,2,2,1,1)}\bigl(J(G)^2\bigr)=1.
\]
This gives a negative answer to a question of Mohammadi on powers of
cover ideals.
\end{abstract}

\maketitle

Let $K$ be a field and $S=K[x_1,\ldots,x_8]$.  For a graph $H$ on
$[8]$, write
\[
I(H)=(x_ix_j:\{i,j\}\in E(H)),
\qquad
J(H)=I(H)^\vee
    =\bigcap_{\{i,j\}\in E(H)}(x_i,x_j).
\]
Thus $J(H)$ is generated by the monomials corresponding to the minimal
vertex covers of $H$.  A graph is Cohen--Macaulay over $K$ when
$S/I(H)$ is Cohen--Macaulay.  We write
\[
\beta_{i,\mathbf b}(M)
  =\dim_K\operatorname{Tor}_i^S(M,K)_{\mathbf b}.
\]

Mohammadi asked whether every ordinary power of the cover ideal of a
Cohen--Macaulay graph again has a linear resolution
\cite[Question~4.1]{Mohammadi}; by the Eagon--Reiner theorem
\cite{EagonReiner}, a graph is Cohen--Macaulay over $K$ precisely when
its cover ideal has a linear resolution over $K$.  H\`a and Van Tuyl
restate the question as \cite[Question~6.2]{HaVanTuyl}:
\begin{quote}
Let $G$ be a graph such that the cover ideal $J(G)$ has a linear
resolution.  Does $J(G)^s$ have a linear resolution for all $s\geq1$?
If not, what hypotheses are needed on $G$ to give this conclusion?
\end{quote}
Mohammadi's positive result
for cactus graphs was later supplied with a corrected proof by Lu and
Wang \cite{LuWang}.  The following example answers the general
question in the negative already for the square.

\begin{theorem}\label{thm:main}
Let $G$ be the graph on $[8]$ with
\[
E(G)=\{14,16,17,23,26,28,37,48,56,57,58\},
\]
where $ij$ denotes the edge $\{i,j\}$.  Then $G$ is Cohen--Macaulay
over every field, $J(G)$ has a $5$-linear resolution, and
\[
\beta_{1,\mathbf a}\bigl(J(G)^2\bigr)=1,
\qquad
\mathbf a=(1,1,2,2,2,2,1,1).
\]
Consequently, $J(G)^2$ does not have a linear resolution over any
field.
\end{theorem}

\begin{proof}
For distinct indices, write $i_1\cdots i_r$ for
$\{i_1,\ldots,i_r\}$.  Splitting the maximal independent sets of $G$
according to their intersection with $\{3,4\}$ gives
\[
\begin{array}{ll}
\text{both }3,4\colon & 345,\ 346;\\
3\ \text{but not}\ 4\colon & 135,\ 138,\ 368;\\
4\ \text{but not}\ 3\colon & 245,\ 247,\ 467;\\
\text{neither}\colon & 125,\ 678.
\end{array}
\]
Hence $\Delta=\Ind(G)$ has exactly $2+3+3+2=10$ facets, each of size
$3$.  Reordered, they are
\[
345,\ 135,\ 245,\ 125,\ 346,\ 247,\ 467,\ 138,\ 368,\ 678.
\tag{1}\label{eq:shelling}
\]
For the second through tenth facets, the facets of the intersection
with the preceding subcomplex are, respectively,
\[
35;\quad 45;\quad 15,25;\quad 34;\quad 24;\quad
46,47;\quad 13;\quad 36,38;\quad 67,68.
\]
Each intersection is nonempty and pure of dimension one.  Hence
\eqref{eq:shelling} is a shelling of the pure two-dimensional complex
$\Delta$.  Thus
\[
K[\Delta]=S/I(G)
\]
is Cohen--Macaulay over every field.  The Eagon--Reiner theorem
\cite{EagonReiner} now implies that $J(G)=I(G)^\vee$ has a linear
resolution.  Since every facet of $\Delta$ has size $3$, every minimal
vertex cover has size $5$, so this resolution is $5$-linear.

Let $C_i$ be the complement in $[8]$ of the $i$th facet in
\eqref{eq:shelling}.  The minimal vertex covers are
\[
\begin{array}{lllll}
C_1=12678, & C_2=24678, & C_3=13678, & C_4=34678, & C_5=12578,\\
C_6=13568, & C_7=12358, & C_8=24567, & C_9=12457, & C_{10}=12345.
\end{array}
\]
Set $x_C=\prod_{i\in C}x_i$ and
\[
m=x^{\mathbf a}
  =x_1x_2x_3^2x_4^2x_5^2x_6^2x_7x_8.
\]
Consider the upper Koszul complex
\[
\mathcal K^{\mathbf a}
 =\left\{
 F\subseteq[8]:
 \frac{m}{x_F}\in J(G)^2
 \right\}.
\]
Since $\mathbf a$ has all coordinates positive, $m/x_F$ is a monomial
for every $F\subseteq[8]$, and since $J(G)^2$ is an ideal,
$m/x_{F'}=(m/x_F)\,x_{F\setminus F'}$ lies in $J(G)^2$ whenever
$F'\subseteq F$ and $F\in\mathcal K^{\mathbf a}$.  Thus
$\mathcal K^{\mathbf a}$ is closed under passing to subsets, hence is a
simplicial complex.

The identities
\[
\frac{m}{x_3x_4}=x_{C_6}x_{C_8},
\qquad
\frac{m}{x_5x_6}=x_{C_4}x_{C_{10}}
\tag{2}\label{eq:positive}
\]
show that $34$ and $56$ are faces.

Since $J(G)^2$ is a monomial ideal, membership is equivalent to
divisibility by $x_{C_i}x_{C_j}$ for some
$1\leq i\leq j\leq10$.  The following table excludes four vertices
and the four cross-edges between $\{3,4\}$ and $\{5,6\}$:
\[
\begin{array}{c@{\qquad}c@{\qquad}c@{\qquad}c}
\toprule
F & \text{condition on each factor cover}
  & \text{eligible covers} & t\\
\midrule
\{1\}\ \text{or}\ \{3,5\}
  & \text{omit }1\ \text{or contain }4,6
  & C_2,C_4,C_8 & 7\\
\{2\}\ \text{or}\ \{4,5\}
  & \text{omit }2\ \text{or contain }3,6
  & C_3,C_4,C_6 & 8\\
\{7\}\ \text{or}\ \{4,6\}
  & \text{omit }7\ \text{or contain }3,5
  & C_6,C_7,C_{10} & 1\\
\{8\}\ \text{or}\ \{3,6\}
  & \text{omit }8\ \text{or contain }4,5
  & C_8,C_9,C_{10} & 2\\
\bottomrule
\end{array}
\tag{3}\label{tab:exclusions}
\]
For a singleton $F=\{r\}$ occurring in
\eqref{tab:exclusions}, the $x_r$-exponent of $m/x_F$ is zero, so both
factor covers must omit $r$.  For a two-element $F$ in the table,
$m/x_F$ and $x_{C_i}x_{C_j}$ both have degree $10$; divisibility would
therefore be equality, and every exponent-two coordinate of $m/x_F$
must occur in both covers.  In each row, the two conditions correspond
respectively to the two choices of $F$.  These observations give the
eligible covers in \eqref{tab:exclusions}.  In each row every eligible
cover contains $t$, so a product of two such covers has
$x_t$-exponent $2$, whereas $m/x_F$ has $x_t$-exponent $1$.  Thus none
of the eight displayed faces belongs to $\mathcal K^{\mathbf a}$.

Because $1,2,7,8$ are not vertices of $\mathcal K^{\mathbf a}$, every
face is contained in $\{3,4,5,6\}$.  Any subset of this set not
contained in $34$ or $56$ contains one of the excluded cross-edges
$35,36,45,46$.  Together with \eqref{eq:positive}, this proves
\[
\mathcal K^{\mathbf a}=\langle34,56\rangle,
\]
the disjoint union of two edges.  The upper-Koszul formula
\cite[Theorem~1.34]{MillerSturmfels} gives
\[
\beta_{1,\mathbf a}\bigl(J(G)^2\bigr)
 =\dim_K\widetilde H_0(\mathcal K^{\mathbf a};K)
 =1.
\]
Finally, $J(G)^2$ is generated in degree $10$, while
$|\mathbf a|=12$.  A $10$-linear resolution would have all first
syzygies in degree $11$, so $J(G)^2$ is not linear.
\end{proof}

\begin{remark}
The vertices $1,6,5,7$ induce a $4$-cycle, so $G$ is not chordal.
Thus Theorem~\ref{thm:main} does not address the Herzog--Hibi--Ohsugi
conjecture that every power of the cover ideal of a chordal graph is
componentwise linear \cite[Conjecture~2.5]{HHO2011}, surveyed in
\cite{HaVanTuyl}.

Ficarra and Moradi constructed analogous failures for squarefree
monomial ideals in every degree $3\leq d\leq n-3$
\cite[Theorem~D]{FicarraMoradi}, but their construction is not a graph
cover ideal.  For $d\geq4$ its generators have a nontrivial common
factor, whereas a graph cover ideal has greatest common divisor $1$
because every vertex extends to a maximal independent set.  For
$d=3$, in their notation,
\[
\{a,c,f,y_1,\ldots,y_{n-6}\}
\]
is a minimal hitting set of the generator supports, so the Alexander
dual has a generator of degree $n-3\geq3$ and is not quadratic.
Theorem~\ref{thm:main} gives the graph-cover-ideal counterexample asked
for by Mohammadi.
\end{remark}

\begin{thebibliography}{99}

\bibitem{EagonReiner}
J.~A. Eagon and V.~Reiner,
\emph{Resolutions of Stanley--Reisner rings and Alexander duality},
J. Pure Appl. Algebra \textbf{130} (1998), no.~3, 265--275.

\bibitem{FicarraMoradi}
A.~Ficarra and S.~Moradi,
\emph{Stanley--Reisner ideals with linear powers},
arXiv:2508.10354, 2025.

\bibitem{HaVanTuyl}
H.~T. H\`a and A.~Van Tuyl,
\emph{Powers of componentwise linear ideals:
the Herzog--Hibi--Ohsugi conjecture and related problems},
Res. Math. Sci. \textbf{9} (2022), Paper No.~22, 26 pp.

\bibitem{HHO2011}
J.~Herzog, T.~Hibi and H.~Ohsugi,
\emph{Powers of componentwise linear ideals},
Combinatorial aspects of commutative algebra and algebraic geometry,
49--60, Abel Symp., vol.~6, Springer, Berlin, 2011.

\bibitem{LuWang}
D.~Lu and Z.~Wang,
\emph{Correction to: Powers of the vertex cover ideals
(Collect. Math. 65 (2014), 169--181)},
Collect. Math. \textbf{76} (2025), 205--209.

\bibitem{MillerSturmfels}
E.~Miller and B.~Sturmfels,
\emph{Combinatorial Commutative Algebra},
Graduate Texts in Mathematics, vol.~227,
Springer, New York, 2005.

\bibitem{Mohammadi}
F.~Mohammadi,
\emph{Powers of the vertex cover ideals},
Collect. Math. \textbf{65} (2014), no.~2, 169--181.

\end{thebibliography}

\end{document}